\chapter{Deep Dive: Computer Vision}
\label{chap:cv}

Computer Vision (CV) adalah bidang AI yang melatih komputer untuk menafsirkan dan memahami dunia visual.

\section{Convolutional Neural Networks (CNN)}

Jika LLM menggunakan Transformer, maka Computer Vision menggunakan **CNN**. CNN bekerja dengan cara memindai gambar menggunakan "filter" kecil untuk mendeteksi fitur seperti garis, sudut, dan bentuk.

\begin{tipbox}
\textbf{Analogi Senter:}
Bayangkan Anda melihat lukisan besar di ruangan gelap dengan senter kecil. Anda menyusuri lukisan itu inci demi inci. CNN melakukan hal yang sama untuk memahami gambar.
\end{tipbox}

\section{Object Detection: YOLO}

YOLO (You Only Look Once) adalah algoritma populer untuk mendeteksi objek secara real-time. Tidak seperti model lama yang melihat gambar berkali-kali, YOLO melihat sekali dan langsung tahu di mana letak semua objek.

\begin{codeblock}[title=Konsep Output YOLO]
Input: Gambar jalan raya
Output:
- Mobil (Confidence: 99%, Box: [x=10, y=20, w=50, h=30])
- Orang (Confidence: 85%, Box: [x=60, y=50, w=10, h=20])
- Lampu Merah (Confidence: 90%, Box: [x=90, y=10, w=5, h=10])
\end{codeblock}

\section{Image Segmentation}

Jika Object Detection membuat kotak di sekitar objek, Segmentation mewarnai setiap pixel.
\begin{itemize}
    \item **Semantic Segmentation:** Semua "mobil" diwarnai merah.
    \item **Instance Segmentation:** Mobil A merah, Mobil B biru.
\end{itemize}

\section{Penerapan di Industri}

\begin{itemize}
    \item **Medis:** Mendeteksi tumor di hasil rontgen.
    \item **Manufaktur:** Mendeteksi cacat produk di jalur perakitan.
    \item **Otomotif:** Mobil otonom membaca rambu lalu lintas.
\end{itemize}
